\chapter{Security \& Liveness Detection}

\section{Threat Model}
Presentation Attacks (PAs) aim to subvert biometric authentication by presenting an artifact to the camera sensor. The SHIELD threat model encompasses:
\begin{itemize}
    \item \textbf{Printed Photo Attacks:} A high-resolution 2D image presented to the lens.
    \item \textbf{Screen Replay Attacks:} A pre-recorded video or photograph displayed on a digital screen (e.g., iPad, smartphone).
    \item \textbf{3D Mask Attacks:} Silicone, paper, or resin masks worn over a human face.
    \item \textbf{Deepfake/Virtual Camera Injections:} Synthetically generated face swaps injected directly into the video feed at the OS level.
\end{itemize}

\section{Defense Mechanisms}

\subsection{Passive Texture Analysis (MiniFASNet)}
Digital screens emit distinct chromatic profiles and moiré patterns due to pixel grids, while printed photos lack specular reflections. MiniFASNet excels at identifying these high-frequency artifacts. 
\\ \textbf{Limitation:} A high-quality 3D silicone mask mimics human skin texture, often bypassing 2D convolutional filters.

\subsection{Physiological Liveness (rPPG)}
To counter 3D masks, SHIELD employs rPPG. Synthetic artifacts, regardless of texture quality, do not possess a cardiovascular system. By strictly filtering and classifying the micro-variations in the green color channel of the facial ROI, SHIELD accurately identifies the lack of a human pulse.
\\ \textbf{Limitation:} The rPPG signal is extremely susceptible to ambient lighting noise and severe head motion.

\subsection{Dynamic Behavioral Tracking}
To defeat static photos and pre-recorded videos, the \texttt{ChallengeService} prompts the user with randomized tasks ("Look Left", "Smile"). The \texttt{BehavioralAnalyzer} tracks the Euler angles and MediaPipe landmarks to verify intent in real-time.

\subsection{Environment Security}
To defend against virtual camera injections (e.g., OBS Virtual Cam), the system enforces Safe Exam Browser (SEB) cryptographic header validation. The FastAPI backend rejects all WebSocket handshakes lacking a valid SEB trust token, ensuring the video feed originates from a secure hardware boundary.

\begin{table}[H]
\centering
\caption{Threat vs Defense Matrix}
\begin{tabularx}{\textwidth}{|l|l|X|}
\hline
\textbf{Attack Vector} & \textbf{Primary Defense Mechanism} & \textbf{Effectiveness} \\
\hline
Printed Photograph & MiniFASNet, Behavioral Challenge & Very High \\
Screen Replay & MiniFASNet, rPPG (No pulse) & Very High \\
3D Silicone Mask & rPPG (No pulse), Behavioral Challenge & High \\
Virtual Camera Swap & SEB Token Validation, HTTP Headers & Absolute (By Design) \\
\hline
\end{tabularx}
\end{table}

\begin{table}[H]
\centering
\caption{Liveness Technique Comparison}
\begin{tabularx}{\textwidth}{|X|X|X|}
\hline
\textbf{Technique} & \textbf{Strengths} & \textbf{Computational Overhead} \\
\hline
Texture Analysis (2D) & Instantaneous decision, highly robust to screens. & Low (CPU Optimized via ONNX) \\
Physiological (rPPG) & Unmatched against 3D masks and high-fidelity prints. & High (Requires sliding window buffering) \\
Behavioral Response & Highly interactive, defeats static replays. & Medium (Requires dense facial landmarking) \\
\hline
\end{tabularx}
\end{table}

\section{Conclusion on System Security}
SHIELD implements a defense-in-depth architecture. By aggregating passive spatial metrics, active temporal behavior, and physiological pulse detection through a unified Fusion Engine, the system mathematically ensures that bypassing one modality (e.g., using a high-res screen to bypass rPPG) triggers a critical failure in another (MiniFASNet detecting screen borders).